% =============================================================================
\section{将我们的论文映射到 WRP 矩阵}
\label{sec:mapping}
% =============================================================================

\Cref{tab:interactions} 将我们的每篇论文映射到 WRP 交互矩阵中，展示它们各自覆盖了哪些跨维度单元。那些仍然稀疏的非对角单元，正对应着 \Cref{sec:vision} 中提出的开放机会。

\begin{table}[t]
\centering
\caption{WRP 交互矩阵。每个单元格列出我们已经覆盖该跨维交互的论文，并附上关键外部参考。稀疏单元表示开放机会。}
\label{tab:interactions}
\small
\renewcommand{\arraystretch}{1.4}
\begin{tabularx}{\textwidth}{>{\bfseries\small}l X X X}
\toprule
& \textbf{Workload} & \textbf{Router} & \textbf{Pool} \\
\midrule
\rowcolor{workloadblue!5}
Workload   & \textit{CDF 原型}~\cite{fleetopt2026};
            轨迹~\cite{patel2024splitwise, servegen2025}
  & AVR~\cite{avr2026}（模态）；
    VCD~\cite{vcd2026}（安全）；
    SafetyL1/L2~\cite{safetyl1_2026, safetyl2_2026}（内容安全）；
    FeedbackDet~\cite{feedbackdet2026}（聊天自适应）；
    HaluGate~\cite{halugate2025}（事实验证）；
    \textit{仍较稀疏，见~\Cref{opp:session-cascade,opp:failure-routing}}
  & FleetOpt~\cite{fleetopt2026}（CDF$\to$资源规模）；
    1/W Law~\cite{onewlaw2026}（能耗）；
    \textit{仍较稀疏，见~\Cref{opp:proactive-scaling}} \\
\rowcolor{routergreen!5}
Router & FleetOpt~\cite{fleetopt2026}（压缩改变 CDF）；
         WhenToReason~\cite{wang2025reason}
  & \vsr{}~\cite{vllmsr2026}（信号）；
    ProbPol~\cite{probpol2026}（冲突）；
    FastRouter~\cite{fastrouter2026}（延迟）；
    mmBERT-Embed~\cite{mmberted2025}（嵌入）；
    OATS~\cite{oats2026}（工具）；
    SemCache~\cite{wang2025cache}（缓存）；
    HaluGate~\cite{halugate2025}（验证）
  & FleetOpt~\cite{fleetopt2026}（联合设计 $\gamma$）；
    \textit{仍较稀疏，见~\Cref{opp:dynamic-boundary,opp:kv-retention}} \\
\rowcolor{poolorange!5}
Pool & FleetSim~\cite{fleetsim2026}（拓扑$\to$成本）；
       1/W~\cite{onewlaw2026}（拓扑$\to$能耗）
  & \textit{仍较稀疏，见~\Cref{opp:pool-feedback}}
  & FleetSim~\cite{fleetsim2026}（A10G/A100/H100）；
    M\'{e}lange~\cite{griggs2024melange}；
    DistServe~\cite{zhong2024distserve} \\
\bottomrule
\end{tabularx}
\end{table}
